% ============================================================
\section{Funtor \texorpdfstring{$I: DLat \hookrightarrow Lat$\ e\ $D: Lat \to DLat$}{I: DLat -> Lat e D: Lat -> DLat}}
\label{sec:functor-lat-dlat}

\subsection{Descrição das Categorias}

\begin{definition}[Categoria $\mathbf{Lat}$ (modelo MATLAB/Octave)]
Um objeto é um reticulado finito $(L, \wedge, \vee)$, ou seja, um conjunto finito com duas operações binárias de encontro (\emph{meet}) e junção (\emph{join}). Em MATLAB, representa-se por uma \texttt{struct} com campos \texttt{elements}, \texttt{meet} e \texttt{join}.
\end{definition}

\begin{lstlisting}[style=matlab, caption={Objeto em $\mathbf{Lat}$}]
L.elements = [0, 1];
L.meet = @(x, y) min(x, y);   % encontro
L.join = @(x, y) max(x, y);   % juncao
\end{lstlisting}

\begin{definition}[Categoria $\mathbf{DLat}$ (modelo MATLAB/Octave)]
Um objeto é um reticulado distributivo finito $(L, \wedge, \vee)$, onde as operações satisfazem adicionalmente a lei distributiva:
\[
    x \wedge (y \vee z) = (x \wedge y) \vee (x \wedge z).
\]
A representação em MATLAB é idêntica à de $\mathbf{Lat}$, acrescentando a verificação da distributividade.
\end{definition}

\begin{lstlisting}[style=matlab, caption={Verificação da distributividade em $\mathbf{DLat}$}]
function ok = is_distributive(L)
    ok = true;
    for x = L.elements
        for y = L.elements
            for z = L.elements
                lhs = L.meet(x, L.join(y, z));
                rhs = L.join(L.meet(x, y), L.meet(x, z));
                if lhs ~= rhs, ok = false; return; end
            end
        end
    end
end
\end{lstlisting}

\subsection{Funtor de Inclusão \texorpdfstring{$I: DLat \hookrightarrow Lat$}{I: DLat -> Lat}}

Existe o funtor de inclusão covariante:
\[
    I: \mathbf{DLat} \hookrightarrow \mathbf{Lat},
\]
que vê cada reticulado distributivo simplesmente como um reticulado, esquecendo a propriedade adicional de distributividade. Nos morfismos, $I$ atua como a identidade: qualquer homomorfismo de reticulados distributivos é também um homomorfismo de reticulados.

\begin{lstlisting}[style=matlab, caption={Funtor de inclusão $I: \mathbf{DLat} \hookrightarrow \mathbf{Lat}$}]
function L = inclusion_DLat_to_Lat(DL)
    % I: esquece a propriedade de distributividade
    L.elements = DL.elements;
    L.meet = DL.meet;
    L.join = DL.join;
end
\end{lstlisting}

\subsection{Reflector Distributivo \texorpdfstring{$D: Lat \to DLat$}{D: Lat -> DLat}}

Existe também um reflector distributivo:
\[
    D: \mathbf{Lat} \to \mathbf{DLat},
\]
que envia um reticulado $L$ no seu maior quociente distributivo
\[
    D(L) = L / \theta_{\mathrm{dist}},
\]
onde $\theta_{\mathrm{dist}}$ é a menor congruência que força a distributividade.

\begin{remark}
A restrição de $\mathbf{Lat}$ aos reticulados distributivos não define um funtor em toda a categoria $\mathbf{Lat}$, porque nem todo o reticulado é distributivo. O que existe naturalmente é a inclusão $I$ ou então o quociente reflector $D$.
\end{remark}

\subsection{Transformações Naturais \texorpdfstring{$q: \mathrm{Id}_{\mathbf{Lat}} \Rightarrow I \circ D$}{q: Id -> I . D}}

Existe uma transformação natural
\[
    q: \mathrm{Id}_{\mathbf{Lat}} \Rightarrow I \circ D,
\]
em que a componente $q_L$ é a projeção canónica de $L$ no seu quociente distributivo $D(L)$:
\[
    q_L: L \twoheadrightarrow L/\theta_{\mathrm{dist}} = D(L).
\]
A naturalidade decorre do facto de que, para qualquer morfismo $f: L \to M$ em $\mathbf{Lat}$, o morfismo $D(f)$ está bem definido no quociente e o diagrama de naturalidade comuta.

\begin{lstlisting}[style=matlab, caption={Projeção canónica $q_L: L \to D(L)$ (esquema simplificado)}]
function [classes, proj] = distributive_quotient(L)
    % Calcula as classes de equivalencia da menor congruencia
    % que forca a distributividade.
    % proj(i) indica o indice da classe do elemento L.elements(i).
    n = length(L.elements);
    proj = 1:n;  % inicialmente, cada elemento e a sua propria classe
    changed = true;
    while changed
        changed = false;
        for x = 1:n
            for y = 1:n
                for z = 1:n
                    lhs = proj(L.meet(L.elements(x), ...
                               L.join(L.elements(y), L.elements(z))) + 1);
                    rhs_l = proj(L.meet(L.elements(x), L.elements(y)) + 1);
                    rhs_r = proj(L.meet(L.elements(x), L.elements(z)) + 1);
                    rhs = proj(L.join(L.elements(rhs_l - 1), ...
                               L.elements(rhs_r - 1)) + 1);
                    if lhs ~= rhs
                        % fundir as duas classes
                        old = max(lhs, rhs);
                        new = min(lhs, rhs);
                        proj(proj == old) = new;
                        changed = true;
                    end
                end
            end
        end
    end
    classes = unique(proj);
end
\end{lstlisting}